\documentclass[12pt,a4paper]{article}

% 基础包
\usepackage{ctex}
\usepackage{geometry}
\geometry{left=2.5cm,right=2.5cm,top=2.5cm,bottom=2.5cm}

% 数学
\usepackage{amsmath}
\usepackage{amssymb}

% 颜色与盒子
\usepackage{xcolor}
\usepackage{tcolorbox}
\tcbuselibrary{skins}

% 取消页码
\pagestyle{empty}

% 自定义颜色
\definecolor{mainblue}{RGB}{70,130,180}
\definecolor{lightgray}{RGB}{245,245,245}

% 段落间距
\setlength{\parskip}{0.8em}

\begin{document}

% 人工智能
\section*{\color{mainblue}\large 人工智能}

\begin{tcolorbox}[
    colback=white,
    colframe=mainblue,
    boxrule=0.5pt,
    arc=2mm,
    left=15pt,right=15pt,top=10pt,bottom=10pt
]
\textbf{任务}\quad
$\begin{cases}
    \text{预测}\\
    \text{决策}
\end{cases}$

\vspace{0.6em}

\textbf{技术}\quad
$\begin{cases}
    \text{搜索}\ \sim\ \text{探索分支好坏}\\
    \text{推理}\ \sim\ \text{旧知识}\rightarrow\text{新知识}\\
    \text{学习}=\mathsf{Machine\ Learning}=\text{系统通过经验提升性能}\\
    \text{博弈}=\mathsf{Reinforcement\ Learning}
\end{cases}$
\end{tcolorbox}

% 机器学习的数学定义
\section*{\color{mainblue}\large 机器学习的数学定义}

\begin{tcolorbox}[
    colback=lightgray,
    colframe=lightgray,
    boxrule=0pt,
    arc=3mm,
    left=12pt,right=12pt,top=12pt,bottom=12pt
]
针对一个预测任务，其数据集是$\mathcal{D}$，对于一个机器学习预测模型$f$，预测任务的性能指标可以通过一个函数$P(\mathcal{D},f)$来表示，那么机器学习的过程就是在一个给定的模型空间$\mathcal{F}$中，寻找可以最大化性能指标的预测模型$f^*$：

\vspace{1em}

\begin{center}
\large
$
\displaystyle f^* = \mathop{\arg\max}_{f \in \mathcal{F}} P(\mathcal{D},f) = \mathsf{ML}(\mathcal{D})
$
\end{center}
\end{tcolorbox}

\end{document}
